% Top margin, right margin, left margin, bottom margin, footnote skip
\usepackage[utf8]{inputenc}
\usepackage{biblatex}
\addbibresource{./references.bib}
% linktocpage shall be added to snippets.
\usepackage{hyperref,theoremref}
\hypersetup{
	colorlinks, 
	linkcolor={red!40!black}, 
	citecolor={blue!50!black},
	urlcolor={blue!80!black},
	linktocpage % Link table of content to the page instead of the title
}

\usepackage{algorithm}
\usepackage{algpseudocode}
\usepackage{blindtext}
\usepackage{titlesec}
\usepackage{amsthm}
\usepackage{thmtools}
\usepackage{amsmath}
\usepackage{amssymb}
\usepackage{graphicx}
\usepackage{titlesec}
\usepackage{xcolor}
\usepackage{multicol}
\usepackage{breqn}
\usepackage{hyperref}
\usepackage{import}
\usepackage{libertinus}           % Load the Libertinus font
\usepackage{mathbbol}
\usepackage{bm}
\usepackage{mathrsfs}
\usepackage[draft]{tikzit}
\usepackage{titlesec}

\setcounter{secnumdepth}{4}
\setcounter{tocdepth}{4}    

\titleformat{\paragraph}
{\normalfont\normalsize\bfseries}{\theparagraph}{1em}{}
\titlespacing*{\paragraph}
{0pt}{3.25ex plus 1ex minus .2ex}{1.5ex plus .2ex}

\usepackage{chngcntr}
\counterwithin*{equation}{section}
\renewcommand{\theequation}{\thesection.\arabic{equation}}

\usepackage{tikz}
\usepackage{tikz-cd}
\usetikzlibrary{arrows.meta, positioning, calc, arrows}
% \usepackage{calrsfs}


\newtheorem{theorem}{Theorem}[section]
\newtheorem{corollary}[theorem]{Corollarium}

\theoremstyle{definition}
\newtheorem{lemma}[theorem]{Lemma}
\newtheorem{definition}[theorem]{Definition}
\newtheorem{problem}{Problem}


\theoremstyle{remark}
\newtheorem{remark}[theorem]{Remark}
\newtheorem{example}[theorem]{Exampli Gratia}
% Proof Environments
\newcommand{\thm}[2]{\begin{theorem}[#1]{}#2\end{theorem}}

% \renewenvironment{proof}{{\bfseries\emph{Demonstratio.}}}{\qed}

\newcommand{\X}{\vec{\mathcal{X}}}
\newcommand{\Z}{\vec{\mathcal{Z}}}



% face and degenerate maps
\renewcommand{\d}{\partial}
\newcommand{\s}{\sigma}

\newcommand{\Ch}{\textit{\textbf{Ch}}}

\newcommand{\T}{\text{Tor}}
\newcommand{\E}{\text{Ext}}
\newcommand{\id}{\text{id}}

\newcommand{\im}{\text{Im}}

\newcommand{\ket}[1]{|#1\rangle}
\newcommand{\bra}[1]{\langle#1|}

\let\H\relax
\DeclareMathOperator{\H}{\text{Hom}}

\renewcommand{\emph}{\textbf}

\input{figures/tizk-style.tikzstyles}
